\documentclass[12pt, a4paper]{article}


% A pretty common set of packages
\usepackage[margin=2.5cm]{geometry}
\usepackage[T1]{fontenc}
\usepackage{graphicx}
\usepackage{amssymb}
\usepackage{amsmath}
\usepackage{color}
\usepackage{booktabs}
\usepackage{multirow}
\usepackage{engord}
\usepackage{soul}
\usepackage{textcomp}
\usepackage{parskip}
\usepackage{setspace}
\usepackage{titlesec}
\usepackage{fancyhdr}
\pagestyle{fancy}
\usepackage[UKenglish]{babel}
\usepackage[UKenglish]{isodate}
\usepackage[skip=2pt,font=footnotesize,justification=centering]{caption}
% \usepackage{natbib}
\usepackage[colorlinks=true, 
    linkcolor=blue,          % color of internal links
    citecolor=blue,        % color of links to bibliography
    filecolor=blue,      % color of file links
    urlcolor=blue]{hyperref}




% Do you prefer Sans Serif fonts?
%\usepackage{sfmath}
%\renewcommand{\familydefault}{\sfdefault} 




% Make some additional useful commands
\newcommand{\ie}{\emph{i.e.}\ }
\newcommand{\eg}{\emph{e.g.}\ }
\newcommand{\etal}{\emph{et al}}
\newcommand{\sub}[1]{$_{\textrm{#1}}$}
\newcommand{\super}[1]{$^{\textrm{#1}}$}
\newcommand{\degC}{$^{\circ}$C}
\newcommand{\wig}{$\sim$}
\newcommand{\ord}[1]{\engordnumber{#1}}
\newcommand{\num}[2]{$#1\,$#2}
\newcommand{\range}[3]{$#1$-$#2\,$#3}
\newcommand{\roughly}[2]{$\sim\!#1\,$#2}
\newcommand{\area}[3]{$#1 \! \times \! #2\,$#3}
\newcommand{\vol}[4]{$#1 \! \times \! #2 \! \times \! #3\,$#4}
\newcommand{\cube}[1]{$#1 \! \times \! #1 \! \times \! #1$}
\newcommand{\figref}[1]{Figure~\ref{#1}}
\newcommand{\eqnref}[1]{Equation~\ref{#1}}
\newcommand{\tableref}[1]{Table~\ref{#1}}
\newcommand{\secref}[1]{Section \ref{#1}}
\newcommand{\XC}{\emph{exchange-correlation}}
\newcommand{\abinit}{\emph{ab initio}}
\newcommand{\Abinit}{\emph{Ab initio}}
\newcommand{\Lonetwo}{L1$_{2}$}
\newcommand{\Dznt}{D0$_{19}$}
\newcommand{\Dtf}{D8$_{5}$}
\newcommand{\Btwo}{B$_{2}$}
\newcommand{\fcc}{\emph{fcc}}
\newcommand{\hcp}{\emph{hcp}}
\newcommand{\bcc}{\emph{bcc}}
\newcommand{\Ang}{{\AA}}
\newcommand{\inverseAng}{{\AA}$^{-1}$}
%\newcommand{\comment}[1]{}
\newcommand{\comment}[1]{\textcolor{red}{[COMMENT: #1]}}
\newcommand{\more}{\textcolor{red}{[MORE]}}
\newcommand{\red}[1]{\textcolor{red}{#1}}





% Change this to modify look of header and footer
\lhead{}
\chead{}
\rhead{}
\lfoot{}
\cfoot{\thepage{}}
\rfoot{}
\renewcommand{\headrulewidth}{0pt}
\renewcommand{\footrulewidth}{0pt}

\begin{document}

\onehalfspacing


\begin{titlepage}

\begin{center}
\includegraphics[width=1in]{figures/bham_crest}

\vspace{0.3in}

\includegraphics[width=3in]{figures/bham_logo}

\vspace{2in}

{\LARGE Transfer Learning for Alzheimer’s Disease Detection: Adapting Video Classification Models for MRI Scans }

\vspace{0.7in}

{\Large Rhys W. Alexander (2458177)}


\vfill{}
Final project report submitted\\ 
in partial fulfilment for the degree of\\
B.SCI. IN ARTIFICIAL INTELLIGENCE AND COMPUTER SCIENCE
\end{center}

\vspace{0.4in}
Date: \today{}     \hfill{} Project supervisor: \\
Word count: X,XXX   \hfill{} Dr Rickson Mesquita
\end{titlepage}






\tableofcontents

\newpage{}










\section{Abstract}
% 250 words

% - Brief overview of the problem and motivation
% - Summary of methodology and main contributions
% - Key results and conclusions
% - Implications and significance

\section{Introduction}
% 1,000 words

% - **Problem statement**
%   - Challenges in Alzheimer's disease diagnosis
%   - Importance of early and accurate detection
%   - Role of neuroimaging in diagnosis
% - **Motivation**
%   - Clinical importance of automating AD detection
%   - Limitations of current diagnostic approaches
%   - Why T1-weighted MRI is particularly valuable (accessibility, non-invasive, etc.)
% - **Research objectives**
%   - Examine transfer learning from video models to 3D MRI analysis
%   - Compare 3D CNN performance to alternative approaches
%   - Identify brain regions contributing to model decisions
% - **Novel contributions**
%   - Application of pre-trained video classification models for MRI analysis
%   - Domain-specific preprocessing pipeline for structural brain MRI
%   - Subject-level validation methodology preventing data leakage
% - **Dissertation roadmap**
%   - Brief outline of subsequent chapters


\section{Literature review}
% 2,000 words 

% - **Alzheimer's Disease and Neuroimaging**

%   - Pathophysiology with emphasis on structural changes
%   - Medical background of AD (emphasize hippocampal atrophy as key biomarker)
%   - Other neuroimaging markers (ventricular enlargement, cortical thinning)
%   - Limitations of visual assessment by radiologists
%   - Current clinical diagnostic practices and their limitations
%   - Role of structural MRI in diagnosis
%   - Advantages of T1-weighted imaging for AD detection

% - **Deep Learning for Medical Image Analysis**

%   - Evolution from traditional ML to deep learning in medical imaging
%   - 2D vs. 3D approaches for volumetric data
%   - Transfer learning in medical imaging context
%   - Challenges in deep learning for medical imaging (data scarcity, interpretability)

% - **3D Deep Learning Architectures**

%   - 3D CNN architectures (ResNet and variants)
%   - Vision transformers for volumetric data
%   - Video classification models and their adaptation to medical data
%   - Performance comparisons from existing literature

% - **MRI Preprocessing for Deep Learning**

%   - Skull stripping methodologies
%   - Registration and normalization approaches
%   - Impact of preprocessing on model performance
%   - Current best practices

% - **Current State of the Art**
%   - Recent advances in automated AD detection
%   - Performance benchmarks and limitations
%   - Gap addressed by this research


\section{Methodology}
% 2,500 words

% - **Data Acquisition and Characteristics**

%   - ADNI dataset description and selection criteria
%   - Patient demographics and diagnostic criteria
%   - MRI acquisition parameters (focusing on T1w MPRAGE)
%   - Data distribution analysis (balance, demographics)

% - **Comprehensive Preprocessing Pipeline**

%   - DICOM to NIfTI conversion
%   - Skull stripping using SynthStrip (justification over alternatives)
%   - Voxel standardization to 1×1×1mm
%   - Cropping and reshaping strategy (128×128×128)
%   - Bias field correction and orientation standardization
%   - Rationale for omitting spatial normalization

% - **Data Splitting Strategy**

%   - Subject-level splitting methodology
%   - Round-robin approach for balanced distribution
%   - Final distribution statistics (subjects and scans per split)
%   - Prevention of data leakage concerns

% - **Data Augmentation**

%   - Augmentation techniques implemented (affine transformations, noise, gamma)
%   - Justification for chosen techniques
%   - Impact on model generalization

% - **Model Architectures**

%   - 3D ResNet, r3d18, architecture details
%   - Transfer learning from Kinetics400 pre-training
%   - Layer freezing strategy with rationale
%   - Alternative architectures explored - MC318
%   - MViT investigation and memory constraint challenges
%   - Parameter counts and computational considerations
%   - Implementation details - PyTorch, Weights \& Biases

% - **Training Framework and Implementation**

%   - PyTorch implementation with Weights \& Biases integration
%   - Hyperparameter selection process
%   - batch size
%   - Early stopping criteria
%   - Loss function (weighted cross-entropy) and optimization strategy
%   - Learning rate scheduling approach
%   - Hardware configuration and constraints
%   - Computational optimizations attempted

% - **Evaluation Methodology**
%   - Classification metrics selection and justification
%   - Validation strategy
%   - Statistical analysis approach
%   - Cross-validation approach

\subsection{Data Acquisition and Characteristics}

The Alzheimer's Disease Neuroimaging Initiative (ADNI) database served as the primary data source for this study. ADNI represents a comprehensive, longitudinal dataset specifically designed for Alzheimer's disease research, offering rigorously standardized MRI acquisitions with corresponding clinical diagnoses and metadata.

\subsubsection{Dataset Selection and Access}

After evaluating potential neuroimaging repositories (including OASIS), ADNI was selected for its comprehensive coverage, standardized acquisition protocols, and expert-validated diagnoses. Access was obtained through a formal application process describing the research objectives. The analysis initially utilized data from ADNI-1, later expanding to incorporate volumes from ADNI-2, ADNI-3, and ADNI-4 to increase sample diversity and size.

\subsubsection{Image Acquisition Parameters}

All selected scans were T1-weighted Magnetization Prepared Rapid Gradient Echo (MPRAGE) sequences, chosen for their optimal gray/white matter contrast which facilitates hippocampal visualization, standardized acquisition parameters across multiple ADNI sites, high signal-to-noise ratio for structural analysis, and sensitivity to hippocampal atrophy, a primary biomarker for AD progression. Additionally, the widespread clinical availability and established role of MPRAGE in AD assessment made it an ideal choice for this study.

The T1w MPRAGE sequences typically featured field strengths of 1.5T or 3T with approximately 1mm$^3$ isotropic resolution. The acquisition matrix was approximately 256$\times$256$\times$170, with TR/TE parameters standardized according to ADNI protocol to ensure consistency across imaging sites.

\subsubsection{Subject Demographics and Diagnostic Criteria}

Subjects were classified into two distinct diagnostic categories: Alzheimer's Disease (AD) and Cognitively Normal (CN). The AD cohort consisted of subjects meeting NINCDS-ADRDA criteria for probable AD, while the CN cohort comprised control subjects without significant cognitive impairment. The original dataset distribution was approximately 33\% AD and 67\% CN cases. After initial testing revealed potential overfitting issues (discussed later), additional scans were incorporated. During this expansion phase, all available new AD scans were included, with CN subjects carefully sampled to achieve a balanced 50/50 diagnostic distribution in the final dataset to optimize model training.

\subsubsection{Data Distribution Analysis}

The final dataset contained 1,300 T1w MRI scans from 408 unique subjects:
\begin{itemize}
    \item AD cohort: 650 scans from 203 subjects
    \item CN cohort: 650 scans from 205 subjects
\end{itemize}

Following subject-level splitting (detailed in Data Splitting Strategy section), the distribution across partitions was:
\begin{itemize}
    \item Training set: 1,023 scans (512 AD, 511 CN) from 248 subjects (133 AD, 115 CN)
    \item Validation set: 139 scans (69 AD, 70 CN) from 80 subjects (35 AD, 45 CN)
    \item Test set: 138 scans (69 AD, 69 CN) from 80 subjects (35 AD, 45 CN)
\end{itemize}

This distribution ensured each partition contained sufficient samples for robust model training and evaluation while maintaining diagnostic balance. The deliberate focus on AD versus CN classification (excluding Mild Cognitive Impairment) reflects the clearer structural changes observable in established AD, particularly the hippocampal atrophy that serves as a primary biomarker for disease progression.

\section{Results}
% 2,000 words

% - **Overall Performance Metrics**

%   - Classification accuracy, precision, recall, F1-score
%   - ROC curves and AUC analysis
%   - Confusion matrices and interpretation
%   - k fold Cross-validation results and stability analysis
%   - Statistical significance testing
%   - Benchmarking against literature results
%   - Bayesian analysis on representative populations

% - **Architectural Comparisons**

%   - 3D ResNet vs. Mixed Convolution performance
%   - Impact of layer freezing strategies
%   - Parameter efficiency analysis

% - **Preprocessing Impact Analysis**

%   - Effect of different preprocessing steps
%   - Importance of crop-and-reshape vs. simple interpolation
%   - Impact of skull stripping quality

% - **Augmentation Effectiveness**

%   - Comparative analysis of different augmentation strategies
%   - Quantitative impact on model performance

% - **Training Dynamics**

%   - Learning curves analysis
%   - Convergence patterns
%   - Overfitting observations and mitigations

% - **Error Analysis**

%   - Patterns in misclassifications
%   - Subject-level vs. scan-level errors
%   - Potential confounding factors

% - **Visual Results**
%   - Key visualizations from Weights \& Biases
%   - Representative case studies
%   - Visualization of model attention/activation maps XAI


\section{Discussion}
% 2,000 words

% - **Interpretation of Results**

%   - Critical analysis of performance metrics
%   - Analysis of 70% accuracy in clinical context
%   - Comparison with human radiologist performance
%   - Significance relative to existing literature
%   - Analysis of false positives and false negatives

% - **Technical Insights**

%   - Effectiveness of transfer learning from video domain
%   - Value of 3D vs. 2D/3D hybrid approaches
%   - Computational efficiency considerations
%   - Memory constraints and their implications

% - **Model Interpretability** (if implemented)

%   - Insights from XAI analysis
%   - Visualization techniques for model attention/activation
%   - Correlation with known AD-affected regions
%   - Clinical relevance of identified features

% - **Clinical Implications**

%   - Potential utility as a diagnostic aid
%   - Integration into existing clinical workflows
%   - Complementary role to other diagnostic measures

% - **Technical Challenges and Solutions**

%   - Memory optimization strategies
%   - Training time challenges on consumer hardware
%   - Data preprocessing optimization
%   - Hardware limitations and workarounds
%   - Data leakage prevention and subject isolation

% - **Limitations**
%   - Dataset representativeness and potential biases
%   - Focus on binary classification (AD vs. CN)
%   - Technical constraints (resolution, model capacity)
%   - Hardware constraints impact on model selection
%   - Need for prospective validation
%   - Generalizability concerns


\section{Conclusions}
% 1,000 words

% - **Summary of Contributions**

%   - Key findings on transfer learning effectiveness
%   - Revisiting research objectives
%   - Technical innovations in preprocessing pipeline
%   - Methodological contributions (subject-level validation)

% - **Future Directions**

%   - Architectural improvements
%   - Multi-class classification (including MCI)
%   - Multimodal approaches
%   - Longitudinal analysis potential
%   - Clinical validation pathway
%   - Integration of additional MRI sequences
%   - Consideration of larger/deeper architectures with more compute

% - **Broader Impact**
%   - Implications for AI in neuroimaging
%   - Potential for improving AD diagnosis workflow
%   - Ethical considerations and responsible deployment




\bibliographystyle{IEEEtran}
\bibliography{references} 

\appendix

% - **Detailed Implementation Specifics**

%   - Code snippets for key components
%   - Hyperparameter configurations
%   - Detailed architectures

% - **Additional Visualizations**

%   - Extended results tables
%   - Additional performance metrics
%   - Sample preprocessing visualizations
%   - Extended XAI visualizations

% - **Computational Resources Analysis**
%   - Detailed training times
%   - Memory usage patterns
%   - Optimization attempts



\end{document}



